% =========================================================================
% SECTION 6: LIMITATIONS AND CONCLUSION
% =========================================================================
\section{Limitations and Conclusion}
\label{sec:limitations_conclusion}

\subsection{System Limitations}
\label{sec:limitations}
Scientific rigor requires explicitly delineating the operational boundaries of the PRISM-L1.0 Android implementation:
\begin{enumerate}
    \item \textbf{Planar Ground Geometry}: Inverse ray projection assumes flat ground ($Y_w = 0$). Ascending or descending stairs, ramps, and dropped curbs violate this constraint, introducing non-linear geometric distortion.
    \item \textbf{Negative Obstacle Detection}: The pipeline cannot detect drop-offs, floor depressions, or descending stairways. PRISM is designed strictly as an elevated obstacle-awareness aid.
    \item \textbf{Ground-Contact Point Occlusion}: Distance recovery relies on observing the obstacle bottom boundary ($v_{\text{bottom}}$). If feet or bases are occluded by foreground furniture or luggage, distance is overestimated.
    \item \textbf{Visual Texture and Lighting}: The luminance Sobel filter relies on spatial edge gradients. Textureless flat surfaces (e.g., bare white walls) and severe under-illumination ($< 20$\,lux) degrade candidate extraction.
    \item \textbf{Absence of Metric Odometry}: While the IMU reliably detects steps and classifies walking versus stationary states, it does not integrate metric displacement. Relative velocities reflect observed optical motion rather than full spatial odometry.
    \item \textbf{Headphone Dependency for Spatial Audio}: Mode~B spatial audio strictly requires stereo over-ear or bone-conduction headphones. Device loudspeaker or mono output eliminates stereo panning cues.
    \item \textbf{Absence of BVI Participant Validation}: The current evaluation was conducted exclusively under controlled engineering scenarios with engineering evaluators. Formal clinical validation with BVI participants measuring real-world navigation autonomy, safety, and cognitive workload has not yet been performed.
    \item \textbf{Outdoor Generalization}: The prototype was evaluated exclusively in indoor corridor environments; outdoor generalization is unvalidated.
\end{enumerate}

\subsection{Safety and Assistive Aid Disclaimer}
\label{sec:safety_disclaimer}
PRISM is an exploratory engineering prototype. It does \emph{not} guarantee collision prevention, does not replace a white cane or guide dog, and is not clinically certified. It is intended strictly as an awareness layer complementing primary mobility aids.

\subsection{Conclusion}
\label{sec:conclusion}
This paper presented the design, implementation, and engineering evaluation of \textbf{PRISM}, an Android-based real-time indoor obstacle detection and multimodal hazard warning system. The complete pipeline executes natively on Android devices (API 26--34), integrating monocular camera perception, geometric ground-plane validation, multi-target Kalman tracking, IMU-based pitch/roll orientation fusion, accelerometer step detection, a 9-state hazard lifecycle machine, and a multimodal presentation layer delivering rate-limited speech, 44.1\,kHz stereo spatial audio, and tactile haptics.

Across an automated suite of 100 Android unit tests, PRISM verified 100\% mathematical and state consistency. In controlled indoor evaluations spanning 17 real-world video runs (2,827 frames), the system achieved 97.5\% overall obstacle recall, 100\% recall for close hazards ($Z < 1.0$\,m), 0.08\,m forward-depth accuracy, 96.3\% step detection precision, and zero duplicate warning utterances. In controlled multisensory A/B testing, adding stereo spatial audio (Mode~B) reduced wrong-direction evasion errors from 6.4\% to 1.2\% and lowered mean reaction time by 120\,ms compared to speech-and-haptic guidance alone.

Future work will conduct IRB-approved formal clinical studies with BVI participants, develop dedicated negative-obstacle detection algorithms, and integrate metric pedestrian dead reckoning.
